\documentclass[11pt,a4paper]{ctexart}
\usepackage[a4paper,margin=1.8cm,headheight=14pt]{geometry}
\usepackage{amsmath,amssymb,mathtools}
\usepackage{booktabs,tabularx,longtable,array}
\usepackage{enumitem}
\usepackage{tikz}
\usetikzlibrary{arrows.meta,positioning,fit,calc}
\usepackage[most]{tcolorbox}
\usepackage{xcolor}
\usepackage{fancyhdr}
\usepackage{hyperref}
\usepackage{microtype}

\setlength{\parindent}{2em}
\setlength{\parskip}{0.35em}
\setlist[itemize]{itemsep=0.2em,topsep=0.3em}
\setlist[enumerate]{itemsep=0.2em,topsep=0.3em}
\hypersetup{colorlinks=true,linkcolor=blue!45!black,urlcolor=blue!50!black}

\definecolor{mainblue}{RGB}{45,86,140}
\definecolor{deepblue}{RGB}{30,62,102}
\definecolor{softblue}{RGB}{238,245,252}
\definecolor{softgreen}{RGB}{238,249,243}
\definecolor{softorange}{RGB}{253,246,233}
\definecolor{softred}{RGB}{253,239,239}
\definecolor{softgray}{RGB}{246,247,249}
\definecolor{deepgreen}{RGB}{35,105,75}
\definecolor{deeporange}{RGB}{165,98,22}
\definecolor{deepred}{RGB}{150,55,55}

\pagestyle{fancy}
\fancyhf{}
\lhead{408 Course Atlas}
\rhead{四科世界模型与跨科接口}
\cfoot{\thepage}
\renewcommand{\headrulewidth}{0.4pt}

\newtcolorbox{corebox}[1][]{colback=softblue,colframe=mainblue,title=#1,fonttitle=\bfseries,boxrule=0.8pt,arc=2mm}
\newtcolorbox{bridgebox}[1][]{colback=softgreen,colframe=deepgreen,title=#1,fonttitle=\bfseries,boxrule=0.8pt,arc=2mm}
\newtcolorbox{exam}[1][]{colback=softorange,colframe=deeporange,title=#1,fonttitle=\bfseries,boxrule=0.8pt,arc=2mm}
\newtcolorbox{warn}[1][]{colback=softred,colframe=deepred,title=#1,fonttitle=\bfseries,boxrule=0.8pt,arc=2mm}
\newtcolorbox{method}[1][]{colback=softgray,colframe=black!45,title=#1,fonttitle=\bfseries,boxrule=0.6pt,arc=2mm}

\newcommand{\kw}[1]{\textbf{#1}}
\newcommand{\code}[1]{\texttt{#1}}

\title{\vspace{-1.2cm}\textbf{408 Course Atlas}\\[0.45em]
\large 四科世界模型、接口与综合路径}
\author{408 计算机综合方法论总册}
\date{}

\begin{document}
\maketitle

\begin{corebox}[母问题：408 到底在研究什么？]
408 不是四本教材的并排目录。它研究的是：\kw{数据怎样被组织，程序语义怎样落到机器，操作系统怎样管理执行与资源，独立机器又怎样通过网络通信。}

可以把观察路径压缩成：
\[
\boxed{
\text{Data}
\to \text{Program}
\to \text{Machine}
\to \text{Operating System}
\to \text{Network}
}
\]
但这条链只是\kw{系统观察路径}，不是一个万能世界模型。四科必须保留自己的对象、状态和推理语言。
\end{corebox}

\section{为什么不能用一种语言强行统一四科}

同一道综合题会跨越多个层级，但不同学科真正关心的变量不同。若强行用“对象--状态--机制”一套抽象吞掉全部细节，很容易得到漂亮但不可操作的总图。

\begin{center}
\begin{tabularx}{\textwidth}{>{\bfseries}p{2.5cm} X X}
\toprule
学科 & 学科母问题 & 首先观察什么 \\
\midrule
数据结构 & 怎样按 Workload 组织关系，使关键操作以可接受代价完成？ & Relation / Representation / Operation / Invariant / Cost \\
计算机组成原理 & 怎样用有限硬件实现 ISA 对软件承诺的程序语义？ & State / Location / Path / Resource / Timing / Commit \\
操作系统 & 怎样在并发、有限资源和异步事件下提供受保护的执行环境？ & Object / Relation / Queue / Event / Mechanism / Policy \\
计算机网络 & 没有共享内存和全局即时状态时，独立机器怎样完成通信？ & Scope / Name / State Owner / Event / Transition / Feedback \\
\bottomrule
\end{tabularx}
\end{center}

\begin{method}[共同语言只用于切换镜头]
四科可以共享一组跨科观察槽：
\[
\boxed{
\text{Object / State / Scope / Mapping / Resource / Event / Invariant / Cost}
}
\]
它们只帮助你判断当前应该切到哪门学科，不替代任何 Subject Atlas 的本体模型。
\end{method}

\section{数据结构：从 Workload 反推表示}

\begin{corebox}[Data Structure Mother Model]
\[
\boxed{
\text{Workload}
\to \text{Logical Relation}
\to \text{Required Operations}
\to \text{Representation}
\to \text{Invariant}
\to \text{Algorithm}
\to \text{Cost}
}
\]
\end{corebox}

数据结构的根本判断不是“这题属于链表还是树”，而是：\kw{当前关系是什么、要支持哪些操作、为了这些操作值得维护什么结构、维护成本是多少。}

复杂度、渐进分析、Cost Vector、Workload-first 与 `Relation $\neq$ Representation` 属于 Atlas Foundation，而不是独立机制 Topic。

核心 Topic 依次覆盖：线性关系与存储表示、栈队列、串与模式匹配、树、Heap、Union-Find、图表示与遍历、图结构算法、查找与有序索引、Hash、内部排序、外部排序。

\section{计算机组成原理：从 ISA 语义到精确提交}

\begin{corebox}[Computer Organization Mother Model]
\[
\boxed{
\text{ISA Semantic}
\to \text{Data Movement}
\to \text{Hardware Path}
\to \text{Timing}
\to \text{Architectural State Commit}
}
\]
\end{corebox}

硬件内部可以缓存、重叠、等待甚至进行局部推测，但最终必须保持软件可见的 ISA 语义。

本学科的八个核心 Topic 是：数据表示与运算、ISA 与机器级程序、CPU 数据通路与控制、流水线、主存硬件、Cache、地址翻译硬件、总线与 I/O 硬件。

\begin{method}[计组起手六问]
\begin{enumerate}
\item 当前软件可见状态是什么？
\item 数据现在在哪里？
\item 它经过哪条硬件路径？
\item 需要哪些共享资源？
\item 数据何时可用，何时锁存？
\item 最终何时改变体系结构状态？
\end{enumerate}
\end{method}

\section{操作系统：在状态变化中维护安全与进展}

\begin{corebox}[Operating System Mother Model]
OS 统一状态可写成：
\[
S=(Objects,Relations,Queues)
\]
推演主线是：
\[
\boxed{
\text{State}
\to \text{Transition}
\to \text{Bad State}
\to \text{Safety / Liveness}
\to \text{Minimal Mechanism}
\to \text{Tradeoff}
}
\]
\end{corebox}

OS Atlas Foundation 负责解释 abstraction、virtualization、protection、user/kernel、程序运行环境等共同入口；真正的机制 Topic 只有五个：执行实体与调度、并发同步与死锁、虚拟内存、I/O、文件系统。

\begin{warn}[不要把 Foundation 伪装成第六个机制 Topic]
基础概念是后续机制共同使用的坐标系。它没有必要为了目录整齐与五个机制 Topic 平级竞争 Ownership。
\end{warn}

\section{计算机网络：没有全局即时状态的通信}

\begin{corebox}[Computer Network Mother Model]
\[
\boxed{
\text{Scope}
\to \text{State}
\to \text{Event}
\to \text{Transition}
\to \text{Feedback}
\to \text{Cost}
}
\]
\end{corebox}

网络中的同一个词只有放回作用域才有意义：一跳、端到端、一个 AS、整个 Internet，分别由不同对象保存状态并作出决策。

八个核心 Topic 是：通信基础与性能、单跳交付、可靠传输、IP 与分组转发、路由、UDP/TCP、拥塞控制、应用层服务语义。

\begin{method}[网络起手八问]
先确认：Scope、通信对象、名字、状态所有者、事件、状态/封装变化、反馈信号、时间/带宽/队列成本。
\end{method}

\section{Bridge：只有真实交接才值得独立建册}

Bridge 不是“两个知识点有联系”。它必须先通过 Validity，再通过 Standalone Promotion。

\begin{bridgebox}[Gate 1：这是不是一个真实接口？]
删除具体题目后，仍应存在稳定的：
\[
\boxed{
A\ \text{output}
\to \text{translation/shared structure}
\to B\ \text{input}
}
\]
如果只是 B 使用了 A 的既有机制，优先记为 Use；如果只是多个模块完成一个具体过程，则归 Integration。
\end{bridgebox}

\begin{bridgebox}[Gate 2：它值得独立建册吗？]
继续检查四项：\kw{Ownership Pressure / Reuse / Current-Scope Relevance / New Inference}。

真连接但没有独立维护必要，不占一个 Canonical Bridge Owner。
\end{bridgebox}

当前三个 Core Cross-Subject Bridge 为：

\begin{enumerate}
\item \kw{X-B01 Privilege / Exception / System Call $\times$ OS Control}：硬件控制转移怎样成为内核接管执行的入口；
\item \kw{X-B02 Hardware Address Translation $\times$ OS Virtual Memory}：OS 建立映射，硬件消费映射，缺失时 fault 回到 OS 修复；
\item \kw{X-B03 Interrupt / DMA $\times$ OS I/O}：设备完成怎样从硬件事件变成内核 completion 与进程 wakeup。
\end{enumerate}

X-B04 Process / Socket $\times$ Transport Endpoint 结构上是真接口，但是否晋升 Core 仍需 408 真题和 Coverage evidence 支撑。

\begin{warn}[Anti-Bridge：几个常见误区]
\begin{itemize}
\item Hardware Cache $\neq$ OS Page Cache；
\item TLB miss $\neq$ Cache miss $\neq$ Page Fault；
\item Routing algorithm uses graph ideas $\neq$ routing protocol is literally a data-structure algorithm；
\item “系统里用了 Queue/Tree/Hash”通常只是 Use，不自动产生 Data Structure $\times$ Systems Bridge。
\end{itemize}
\end{warn}

\section{Integration：完整过程怎样组合成熟模块}

Integration 拥有具体协作轨迹，不重新定义局部机制。

\begin{bridgebox}[X-I01：一次 LOAD / Memory Access 的完整慢路径]
\[
\text{Instruction}
\to \text{Effective Address}
\to \text{VA}
\to \text{TLB/Page Table}
\to \text{PA}
\to \text{Cache}
\to \text{Memory}
\]
如果发生 Page Fault：
\[
\text{Fault}
\to \text{Kernel}
\to \text{VM Repair}
\to \text{PTE Update}
\to \text{Retry}
\to \text{Commit}
\]
\end{bridgebox}

\begin{bridgebox}[X-I02：一次 Blocking File read 的完整生命周期]
\[
\text{User Process}
\to \text{System Call}
\to \text{fd/OFD/inode}
\to \text{Page Cache / I/O Request}
\to \text{Block}
\to \text{DMA/Interrupt}
\to \text{Completion}
\to \text{Wakeup}
\to \text{Return}
\]
\end{bridgebox}

\section{做题控制：先切对学科，再调机制}

\begin{exam}[综合题的第一动作]
不要先问“这是哪一章”。先问：\kw{题目现在要求我追踪的是关系、硬件路径、OS 状态，还是分布式通信状态？}

一旦确定当前层：
\begin{itemize}
\item 数据结构：Relation $\to$ Representation $\to$ Operation $\to$ Invariant $\to$ Cost；
\item 计组：State $\to$ Location $\to$ Path $\to$ Timing $\to$ Commit；
\item OS：Object/Queue $\to$ Event $\to$ Mechanism/Policy $\to$ New State；
\item 网络：Scope/Name/Owner $\to$ Event $\to$ Transition/Feedback。
\end{itemize}
\end{exam}

\section{怎样使用这本总册}

\begin{enumerate}
\item 第一次学习某一科时，只把本册当地图，不用先背全部跨科关系；
\item 进入对应 Subject Atlas，建立那门学科自己的世界模型；
\item 学到某个 Topic 时，先掌握局部机制；
\item 只有两个 Owner 真正交换状态、数据或控制权时，打开 Bridge；
\item 只有要追踪一个完整真实过程时，打开 Integration；
\item 做完题以后，把失误归到模型、识别、路径、执行/检查/表达或考试决策，而不是无限扩张知识分类。
\end{enumerate}

\begin{corebox}[一页压缩]
408 的主干不是“记住四科所有章节”，而是：
\[
\boxed{
\text{Separate Subject Models}
+ \text{Explicit Interfaces}
+ \text{Canonical Processes}
+ \text{Executable Control}
}
\]

四科先各自讲清，接口只在真实交接处出现，综合只追踪真实过程。这样既能覆盖考试，又不会为了“统一”制造第二套模糊知识。
\end{corebox}

\end{document}
